\documentclass[11pt,a4paper]{article}

% ── Codificación, fuentes y lenguaje ───────────────────────────────────────────
\usepackage[utf8]{inputenc}
\usepackage[T1]{fontenc}
\usepackage[default,scale=0.95]{sourcesanspro}
\renewcommand{\familydefault}{\sfdefault}
\usepackage[spanish,es-noshorthands,es-tabla]{babel}

% ── Geometría y espaciado ─────────────────────────────────────────────────────
\usepackage[top=2.2cm, bottom=2.5cm, left=2.2cm, right=2.2cm, headheight=15pt]{geometry}
\usepackage{setspace}
\setstretch{1.18}
\usepackage{parskip}

% ── Matemáticas y unidades ────────────────────────────────────────────────────
\usepackage{amsmath}
\usepackage{amssymb}
\usepackage{siunitx}

% ── Circuitos ─────────────────────────────────────────────────────────────────
\usepackage[european, straightvoltages]{circuitikz}

% ── Tablas ────────────────────────────────────────────────────────────────────
\usepackage{booktabs}
\usepackage{tabularx}
\usepackage{array}
\usepackage{longtable}
\usepackage{needspace}
\usepackage{colortbl}
\usepackage{enumitem}

% ── Colores, cajas e iconos ───────────────────────────────────────────────────
\usepackage[dvipsnames]{xcolor}
\usepackage{tcolorbox}
\tcbuselibrary{skins, breakable}
\usepackage{fontawesome5}
\usepackage{graphicx}

% ── Listados de código (dark-mode infográfico) ────────────────────────────────
%   Estilo base: fondo oscuro con números de línea. Los bloques lstlisting se
%   envuelven automáticamente en una caja tcolorbox redondeada.
\usepackage{listings}

% ── Secciones con barra de color (infografía) ────────────────────────────────
\usepackage{titlesec}
\titleformat{\section}
  {\normalfont\LARGE\bfseries\color{azulNoche}}
  {\colorbox{cyanNeon}{\color{fondoOscuro}\thesection}}{0.7em}{}
  [\vspace{2pt}\color{cyanNeon}\rule{\linewidth}{1.8pt}]
\titlespacing*{\section}{0pt}{24pt}{10pt}
\titleformat{\subsection}
  {\normalfont\large\bfseries\color{azulNoche}}
  {\textcolor{cyanNeon}{\thesubsection}}{0.7em}{}
  [\vspace{1pt}\color{grisLinea}\rule{\linewidth}{0.6pt}]
\titlespacing*{\subsection}{0pt}{14pt}{6pt}
\titleformat{\subsubsection}
  {\normalfont\normalsize\bfseries\color{azulNoche}}
  {\textcolor{cyanNeon}{\thesubsubsection}}{0.7em}{}
\titlespacing*{\subsubsection}{0pt}{10pt}{4pt}

% ── Cabeceras y pies de página ────────────────────────────────────────────────
\usepackage{fancyhdr}
\usepackage{lastpage}
\pagestyle{fancy}
\fancyhf{}
\renewcommand{\headrulewidth}{0pt}
\renewcommand{\footrulewidth}{0pt}
\fancyfoot[C]{%
  \begin{minipage}{\linewidth}
    \vspace{2pt}
    \color{cyanNeon}\rule{\linewidth}{1.2pt}\\[2pt]
    \centering\color{grisTexto}\footnotesize
    Página \thepage\ de \pageref{LastPage}
  \end{minipage}%
}

% ── Hiperenlaces ──────────────────────────────────────────────────────────────
\usepackage[colorlinks=true, linkcolor=azulNoche, urlcolor=cyanNeon]{hyperref}
\sloppy
\emergencystretch 3em

% ── Paleta de colores — tecnológico dark-mode ──────────────────────────────────
\definecolor{fondoOscuro}{HTML}{0D1B2A}%             Dark navy — fondo banda título
\definecolor{verdeTurquesa}{HTML}{004D40}%           Verde turquesa — bandas de marca
\definecolor{azulNoche}{HTML}{1B3A6B}%               Midnight blue — secciones, marcos
\definecolor{azulMedio}{HTML}{2A5298}%               Azul medio — variante de acento
\definecolor{cyanNeon}{HTML}{00C8E0}%                Cyan eléctrico — acentos primarios
\definecolor{verdeSignal}{HTML}{00B887}%             Verde señal — fortalezas, éxito
\definecolor{naranjaVivo}{HTML}{FF7043}%             Naranja vivo — mejoras, advertencias
\definecolor{rojoAlerta}{HTML}{E53935}%              Rojo alerta — errores
\definecolor{grisPapel}{HTML}{F4F6FA}%               Fondo tarjetas claras
\definecolor{grisLinea}{HTML}{C8D0E0}%               Bordes sutiles
\definecolor{grisTexto}{HTML}{6B7A99}%               Texto secundario
\definecolor{amarilloNota}{HTML}{FFB300}%            Notas / advertencias doradas
\definecolor{codigoFondo}{HTML}{1E2D3D}%             Fondo del bloque de código
\definecolor{codigoComentario}{HTML}{7A8FA6}%        Comentarios
\definecolor{codigoCadena}{HTML}{FF9D5C}%            Cadenas / strings
\definecolor{codigoKeyword}{HTML}{00C8E0}%           Palabras clave
\definecolor{codigoNumero}{HTML}{00E5A0}%            Números

% ── Banda de título: portada infográfica ──────────────────────────────────────
%   Diseño en 2 capas: banda de color (por defecto fondo oscuro) + franja de
%   acento cyan abajo. Uso: \bandaTitulo[color]{título}{subtítulo}.
%   El icono se puede cambiar con \renewcommand{\iconoBanda}{...} (FontAwesome).
\newcommand{\iconoBanda}{microchip}
\newcommand{\bandaTitulo}[3][fondoOscuro]{%
  \begin{tcolorbox}[
    enhanced,
    colback=#1, colframe=#1,
    boxrule=0pt, arc=8pt, outer arc=8pt,
    width=\linewidth,
    top=18pt, bottom=6pt, left=14pt, right=14pt,
    borderline south={3.5pt}{0pt}{cyanNeon},
  ]
    \begin{minipage}[c]{0.07\linewidth}
      \centering
      \textcolor{cyanNeon}{\fontsize{30}{30}\selectfont\faIcon{\iconoBanda}}
    \end{minipage}%
    \hspace{8pt}%
    \begin{minipage}[c]{0.88\linewidth}
      {\fontsize{20}{24}\selectfont\bfseries\color{white}#2}\par\vspace{5pt}%
      \textcolor{cyanNeon!80!white}{\small\faIcon{angle-right}~#3}%
    \end{minipage}
    \vspace{4pt}
  \end{tcolorbox}%
  \vspace{8pt}%
}

% ── Tarjetas de datos (infografía) ────────────────────────────────────────────
\tcbset{
  tarjetaDato/.style={
    enhanced, breakable,
    arc=6pt, outer arc=6pt,
    colback=grisPapel, colframe=azulNoche,
    boxrule=1pt,
    fonttitle=\bfseries\small, coltitle=white,
    attach boxed title to top left={yshift=-3mm, xshift=6mm},
    boxed title style={
      colback=azulNoche, colframe=azulNoche,
      arc=4pt, boxrule=0pt,
    },
    drop fuzzy shadow=azulNoche!30!white,
    top=8pt, bottom=8pt, left=10pt, right=10pt,
  },
  tarjetaIcono/.style={
    enhanced, breakable,
    arc=6pt, outer arc=6pt,
    colback=white, colframe=grisLinea, boxrule=0.8pt,
    fonttitle=\bfseries\small, coltitle=azulNoche,
    borderline west={3pt}{0pt}{cyanNeon},
    drop fuzzy shadow=azulNoche!25!white,
    top=6pt, bottom=6pt, left=10pt, right=10pt,
  },
  cajaTitulo/.style={
    enhanced, breakable,
    arc=6pt, outer arc=6pt,
    colback=azulNoche!10!grisPapel, colframe=azulNoche,
    boxrule=1pt,
    fonttitle=\bfseries, coltitle=white,
    attach boxed title to top left={yshift=-3mm, xshift=6mm},
    boxed title style={
      colback=azulNoche, colframe=azulNoche,
      arc=4pt, boxrule=0pt,
    },
    drop fuzzy shadow=azulNoche!25!white,
    top=8pt, bottom=8pt, left=10pt, right=10pt,
  },
  cajaContenido/.style={
    enhanced, breakable,
    arc=5pt, outer arc=5pt,
    colback=grisPapel, colframe=grisLinea,
    boxrule=0.8pt,
    fonttitle=\bfseries\small, coltitle=azulNoche,
    top=6pt, bottom=6pt, left=10pt, right=10pt,
  },
  cajaFortaleza/.style={
    enhanced, breakable,
    arc=6pt, outer arc=6pt,
    colback=verdeSignal!8!white, colframe=verdeSignal,
    boxrule=1pt,
    borderline west={4pt}{0pt}{verdeSignal},
    fonttitle=\bfseries\small, coltitle=verdeSignal!70!black,
    drop fuzzy shadow=verdeSignal!20!white,
    top=6pt, bottom=6pt, left=10pt, right=10pt,
  },
  cajaMejora/.style={
    enhanced, breakable,
    arc=6pt, outer arc=6pt,
    colback=naranjaVivo!8!white, colframe=naranjaVivo,
    boxrule=1pt,
    borderline west={4pt}{0pt}{naranjaVivo},
    fonttitle=\bfseries\small, coltitle=naranjaVivo!80!black,
    drop fuzzy shadow=naranjaVivo!20!white,
    top=6pt, bottom=6pt, left=10pt, right=10pt,
  },
  cajaRecomendacion/.style={
    enhanced, breakable,
    arc=6pt, outer arc=6pt,
    colback=cyanNeon!6!white, colframe=cyanNeon!60!azulNoche,
    boxrule=1pt,
    borderline west={4pt}{0pt}{cyanNeon},
    fonttitle=\bfseries\small, coltitle=azulNoche,
    drop fuzzy shadow=azulNoche!20!white,
    top=6pt, bottom=6pt, left=10pt, right=10pt,
  },
}

% ── Ícono + texto (encabezados de sección y elementos) ────────────────────────
\newcommand{\iconotexto}[2]{\textcolor{cyanNeon}{\faIcon{#1}}~#2}

% ── Listados de código: estilo dark + auto-envoltura ─────────────────────────
\lstdefinestyle{estiloCodigo}{
    backgroundcolor=\color{codigoFondo},
    basicstyle=\ttfamily\small\color{white},
    commentstyle=\color{codigoComentario}\itshape,
    stringstyle=\color{codigoCadena},
    keywordstyle=\color{codigoKeyword}\bfseries,
    numberstyle=\tiny\color{codigoComentario},
    numbers=left,
    numbersep=10pt,
    showstringspaces=false,
    breaklines=true,
    breakatwhitespace=false,
    tabsize=4,
    frame=none,
    captionpos=b,
    aboveskip=6pt,
    belowskip=4pt,
    xleftmargin=14pt,
    framexleftmargin=14pt,
    literate=
      {á}{{\'a}}1 {é}{{\'e}}1 {í}{{\'i}}1 {ó}{{\'o}}1 {ú}{{\'u}}1
      {Á}{{\'A}}1 {É}{{\'E}}1 {Í}{{\'I}}1 {Ó}{{\'O}}1 {Ú}{{\'U}}1
      {ñ}{{\~n}}1 {Ñ}{{\~N}}1 {ü}{{\"u}}1 {Ü}{{\"U}}1
      {¿}{{?`}}1 {¡}{{!`}}1
}
\lstdefinestyle{estiloPython}{
    style=estiloCodigo,
    language=Python,
}
\lstset{style=estiloCodigo}
\BeforeBeginEnvironment{lstlisting}{%
  \begin{tcolorbox}[
    enhanced, breakable,
    arc=6pt, outer arc=6pt,
    colback=codigoFondo, colframe=azulNoche,
    boxrule=1pt,
    drop fuzzy shadow=azulNoche!25!white,
    top=2pt, bottom=2pt, left=0pt, right=0pt,
  ]%
}
\AfterEndEnvironment{lstlisting}{\end{tcolorbox}}

% ── Cajas de contenido temáticas ──────────────────────────────────────────────
\newtcolorbox{cajaRecuerda}{
    enhanced, breakable,
    arc=6pt, outer arc=6pt,
    colback=azulNoche!10!grisPapel,
    colframe=azulNoche, boxrule=1pt,
    borderline west={4pt}{0pt}{cyanNeon},
    fonttitle=\bfseries\small\color{white},
    title={\faIcon{info-circle}~Recuerda},
    attach boxed title to top left={yshift=-3mm, xshift=6mm},
    boxed title style={colback=azulNoche, arc=4pt, boxrule=0pt},
    drop fuzzy shadow=azulNoche!25!white,
    left=10pt, right=10pt, top=8pt, bottom=8pt,
}

\newtcolorbox{cajaConcepto}{
    enhanced, breakable,
    arc=6pt, outer arc=6pt,
    colback=verdeSignal!8!white,
    colframe=verdeSignal, boxrule=1pt,
    borderline west={4pt}{0pt}{verdeSignal},
    fonttitle=\bfseries\small\color{white},
    title={\faIcon{lightbulb}~Concepto clave},
    attach boxed title to top left={yshift=-3mm, xshift=6mm},
    boxed title style={colback=verdeSignal!80!black, arc=4pt, boxrule=0pt},
    drop fuzzy shadow=verdeSignal!20!white,
    left=10pt, right=10pt, top=8pt, bottom=8pt,
}

\newtcolorbox{cajaEjemplo}{
    enhanced, breakable,
    arc=6pt, outer arc=6pt,
    colback=cyanNeon!5!white,
    colframe=cyanNeon!70!azulNoche, boxrule=1pt,
    borderline west={4pt}{0pt}{cyanNeon},
    fonttitle=\bfseries\small\color{fondoOscuro},
    title={\faIcon{code}~Ejemplo},
    attach boxed title to top left={yshift=-3mm, xshift=6mm},
    boxed title style={colback=cyanNeon, arc=4pt, boxrule=0pt},
    drop fuzzy shadow=azulNoche!20!white,
    left=10pt, right=10pt, top=8pt, bottom=8pt,
}

\newtcolorbox{cajaEjercicio}{
    enhanced, breakable,
    arc=6pt, outer arc=6pt,
    colback=azulMedio!7!white,
    colframe=azulMedio, boxrule=1pt,
    borderline west={4pt}{0pt}{azulMedio},
    fonttitle=\bfseries\small\color{white},
    title={\faIcon{pencil-alt}~Ejercicio},
    attach boxed title to top left={yshift=-3mm, xshift=6mm},
    boxed title style={colback=azulMedio, arc=4pt, boxrule=0pt},
    drop fuzzy shadow=azulNoche!20!white,
    left=10pt, right=10pt, top=8pt, bottom=8pt,
}

\newtcolorbox{cajaReto}{
    enhanced, breakable,
    arc=6pt, outer arc=6pt,
    colback=naranjaVivo!6!white,
    colframe=naranjaVivo, boxrule=1pt,
    borderline west={4pt}{0pt}{naranjaVivo},
    fonttitle=\bfseries\small\color{white},
    title={\faIcon{fire}~Reto},
    attach boxed title to top left={yshift=-3mm, xshift=6mm},
    boxed title style={colback=naranjaVivo, arc=4pt, boxrule=0pt},
    drop fuzzy shadow=naranjaVivo!20!white,
    left=10pt, right=10pt, top=8pt, bottom=8pt,
}

% ── Indicadores de dificultad ─────────────────────────────────────────────────
\newcommand{\facil}{\textcolor{verdeSignal}{\faIcon{circle}\,\textbf{Fácil}}}
\newcommand{\intermedio}{\textcolor{naranjaVivo}{\faIcon{circle}\,\textbf{Intermedio}}}
\newcommand{\dificil}{\textcolor{rojoAlerta}{\faIcon{circle}\,\textbf{Difícil}}}

% ─────────────────────────────────────────────────────────────────────────────
\renewcommand{\iconoBanda}{shield-alt}
\fancyhead[L]{\footnotesize\color{grisTexto}\faIcon{shield-alt}~CYBERSEGURIDAD}
\fancyhead[R]{\footnotesize\color{grisTexto}Semana 02}
\hypersetup{pdftitle={Ciberseguridad Aplicada},pdfauthor={Prof. César Rodríguez}
  pdfsubject={CYBERSEGURIDAD}}

\begin{document}
\bandaTitulo[fondoOscuro]{Semana 02 --- Diagnosticar TLS y elegir parametros con medicion}{Grupo 2 $\cdot$ rol Azul $\cdot$ ejemplo de trabajo}

\begin{center}\small\color{grisTexto}Prof. César Rodríguez \textbar{} ELECTRONICA \textbar{} 1 de noviembre de 2026\end{center}

\begin{tabularx}{\linewidth}{@{}>{\raggedright\arraybackslash}p{4.6cm}X@{}}
\toprule
\textbf{Curso} & CYBERSEGURIDAD --- Ciberseguridad Aplicada \\
\textbf{Semana} & 02 de 10 \\
\textbf{Fase} & 1 --- Fundamentos, criptografía y marco legal \\
\textbf{Grupo} & 2 \\
\textbf{Rol} & Azul \\
\textbf{Entrega} & GitHub, carpeta semana\_02/ \\
\textbf{Corte} & Domingo 23:59 de la semana correspondiente; lo posterior no se corrige. \\
\bottomrule
\end{tabularx}

\needspace{5\baselineskip}
\section{\iconotexto{user-shield}{Tu rol en la práctica}}

Parte de su propia captura cruda y aplica el control equivalente: configuración, parche o regla de red. Repite el ataque para demostrar que falla y documenta el antes y el después.

\begin{itemize}[leftmargin=1.4em,itemsep=2pt,topsep=3pt]
  \item Leer su propia evidencia cruda y elegir el control que corta el ataque en el punto más barato.
  \item Aplicar el control y repetir el ataque para demostrar que ya no funciona.
  \item Documentar el antes y el después con los mismos comandos y la diferencia medida.
\end{itemize}

\needspace{5\baselineskip}
\section{\iconotexto{flag}{Qué aporta este rol al entregable}}

La mitigación aplicada, el comando de verificación y la comparación de resultados.

\needspace{5\baselineskip}
\section{\iconotexto{list-ol}{Ejemplo paso a paso}}

El rol Azul parte de una configuracion que ya existe y comprueba si valida: lee el certificado, la cadena y el algoritmo, y elige los parametros de derivacion con una medicion propia.

\subsection{Fijar la linea base del certificado y la cadena}

Se inspecciona un sitio real y se guarda la salida completa. Emisor, fechas y algoritmo se leen del certificado, no de la pagina web.

\begin{lstlisting}[style=estiloCodigo]
mkdir -p semana_02/evidencia
openssl s_client -connect www.example.com:443 -servername www.example.com \
  -verify_return_error </dev/null > semana_02/evidencia/tls_handshake.txt 2>&1
openssl s_client -connect www.example.com:443 -servername www.example.com \
  -showcerts </dev/null 2>/dev/null \
  | openssl x509 -noout -issuer -subject -dates \
  | tee semana_02/evidencia/cert_resumen.txt
\end{lstlisting}

\subsection{Leer algoritmo, clave publica y verificacion de la cadena}

La cadena se valida contra el almacen del sistema. Si falla la validacion, esa linea del registro es el hallazgo: no se cambia el comando para que pase.

\begin{lstlisting}[style=estiloCodigo]
openssl s_client -connect www.example.com:443 -servername www.example.com \
  -showcerts </dev/null 2>/dev/null | openssl x509 -noout -text \
  | grep -E 'Signature Algorithm|Public-Key|DNS:' \
  | tee semana_02/evidencia/cert_detalle.txt
openssl s_client -connect www.example.com:443 -servername www.example.com \
  -CAfile /etc/ssl/certs/ca-certificates.crt -verify_return_error </dev/null \
  >/dev/null 2>semana_02/evidencia/cadena.txt; \
  grep -E 'Verify return code' semana_02/evidencia/cadena.txt
\end{lstlisting}

\subsection{Comparar PBKDF2 y Argon2 con el mismo criterio}

Se mide el coste de cada opcion con el mismo hardware y se elige la que sostiene el mismo tiempo objetivo con mas resistencia. Si Argon2 no esta instalado, se documenta la limitacion y se decide solo con PBKDF2.

\begin{lstlisting}[style=estiloCodigo]
python3 - <<'PY' | tee semana_02/evidencia/kdf_comparativa.txt
import hashlib, time
for iteraciones in (100_000, 600_000):
    t = time.perf_counter()
    hashlib.pbkdf2_hmac('sha256', b'clave', b'sal', iteraciones)
    print('pbkdf2', iteraciones, round(time.perf_counter() - t, 3), 's')
try:
    from argon2.low_level import hash_secret_raw, Type
    t = time.perf_counter()
    hash_secret_raw(b'clave', b'sal', time_cost=3, memory_cost=65536,
                    parallelism=4, hash_len=32, type=Type.ID)
    print('argon2id', round(time.perf_counter() - t, 3), 's')
except ImportError:
    print('argon2 no disponible: se documenta como limitacion')
PY
\end{lstlisting}

\subsection{Verificar que el repositorio no filtra secretos}

Antes de entregar se busca material sensible con patrones y con el historial de git. Un secreto en un commit viejo sigue siendo un secreto expuesto.

\begin{lstlisting}[style=estiloCodigo]
git grep -nE 'BEGIN (RSA|OPENSSH|EC|PRIVATE) PRIVATE KEY' || true
git grep -nEi '(api[_-]?key|secret|token|password)[[:space:]]*=' || true
git log --all --oneline -- semana_02/evidencia/*.pem || true
\end{lstlisting}

\needspace{8\baselineskip}\subsection{Configuracion observada y decision}

\begin{center}
\begin{minipage}{\linewidth}
\centering
{\footnotesize
\textbf{Configuracion observada y decision}\par\smallskip
\begin{tabular}{@{}>{\raggedright\arraybackslash}p{2.40cm}>{\raggedright\arraybackslash}p{3.20cm}>{\raggedright\arraybackslash}p{2.40cm}>{\raggedright\arraybackslash}p{3.20cm}>{\raggedright\arraybackslash}p{2.90cm}@{}}
\toprule
\textbf{Elemento} & \textbf{Estado observado} & \textbf{Riesgo} & \textbf{Decision} & \textbf{Verificacion} \\
\midrule
Certificado & Emisor, fechas y DNS leidos & Caducidad o nombre incorrecto & Se documenta y se vigila la fecha & Salida en cert\_resumen.txt \\
Cadena & Verify return code 0 contra el almacen del sistema & Cadena incompleta & Sin accion si valida & Salida en cadena.txt \\
Algoritmo de firma & Leido del certificado & Algoritmo debil o clave corta & Se reporta si es menor de 2048 bits & Linea en cert\_detalle.txt \\
KDF & Tiempo medido por iteraciones & Coste insuficiente & 600000 iteraciones como minimo & Tabla de tiempos \\
Secretos & Busqueda en arbol e historial & Clave privada versionada & Se retira y se rota & git grep sin resultados \\
\bottomrule
\end{tabular}
}
\par\smallskip{\footnotesize La decision se justifica con la salida pegada; "correcto" sin evidencia no se acepta como verificacion.}
\end{minipage}
\end{center}

\needspace{5\baselineskip}
\section{\iconotexto{box-open}{Qué guardar en el repositorio}}

El informe incluye la lectura del certificado y su cadena, la comparacion de KDF con tiempos propios y la lista de verificaciones ejecutadas antes de entregar.

\begin{lstlisting}[style=estiloCodigo]
semana_02/
|-- informe.md
|-- evidencia/
|   |-- tls_handshake.txt
|   |-- cert_resumen.txt
|   |-- cert_detalle.txt
|   |-- cadena.txt
|   |-- kdf_comparativa.txt
|   `-- secretos.txt
`-- README.md
\end{lstlisting}

\begin{itemize}[leftmargin=1.4em,itemsep=2pt,topsep=3pt]
  \item La validacion de la cadena se ejecuta con el almacen del sistema, no con uno propio.
  \item La eleccion de iteraciones se apoya en tiempos medidos y repetidos.
  \item La busqueda de secretos cubre el historial, no solo el arbol actual.
  \item Si Argon2 no esta disponible, la limitacion esta escrita.
\end{itemize}

\needspace{5\baselineskip}
\section{\iconotexto{ruler-combined}{Cómo se corrige}}

\begin{tabularx}{\linewidth}{@{}X r@{}}
\toprule
Criterio & Peso\\
\midrule
Cifrado y descifrado correctos con verificación de integridad & 3,0 \\
Firma y verificación demostradas, incluido el caso alterado & 2,5 \\
Lectura correcta del certificado y de la cadena & 2,0 \\
Medición de KDF con conclusión sobre coste y seguridad & 1,5 \\
Explicación de qué propiedad cubre cada primitiva & 1,0 \\
\midrule
\textbf{Total} & \textbf{10,0}\\
\bottomrule
\end{tabularx}

\begin{itemize}[leftmargin=1.4em,itemsep=2pt,topsep=3pt]
  \item No se pega la clave privada ni el contenido del almacen de confianza en el informe.
  \item \{'Un certificado valido no implica que el servicio sea seguro': 'se separan ambos juicios.'\}
\end{itemize}

\begin{cajaRecuerda}El trabajo es independiente: cada grupo usa su propio repositorio y su propio laboratorio, y la comparación de resultados es en clase. Corte: Domingo 23:59 de la semana correspondiente; lo posterior no se corrige.
\end{cajaRecuerda}
\end{document}
